% GENERATED FILE. Edit canonical lessons, then run npm run generate:books.
% canonical-source-hash: fnv1a64:98399bf2700e3989
% canonical-lessons: ZH-C371-shenzhi, ZH-C371-suiran, ZH-W371-sui-though, ZH-C371-jielun, ZH-C371-zongjie, ZH-C371-yuanze

\chapter{Although}
\label{ch:zh-although}
\hlchaptermodality{371}

\begin{chapteropening}
\textbf{By the end of this chapter:} \emph{I can say although, go so far as to say even, and draw a conclusion.}

Complete the last lesson of chapter 371: I can say although, go so far as to say even, and draw a conclusion.
\end{chapteropening}

\section[shènzhì]{shènzhì — even (going so far as)}

\label{lesson:ZH-C371-shenzhi}



\begin{quote}
Before the new one: say the Mandarin for the effect something has, then the Mandarin for in short.
\end{quote}

\subsection*{You'll want to know: shènzhì}
\emph{shènzhì} — \textquotedblleft{}even (going so far as)\textquotedblright{}.

\begin{quote}
\emph{tā hěn máng shènzhì méiyǒu shíjiān chī fàn} — He's so busy he hasn't even got time to eat.
\end{quote}

Its characters come later; for now it is a word for the ear.

\begin{cousinweb}
\emph{shènzhì} is built from \emph{shèn} and \emph{zhì}, both inherited from Middle Chinese, \emph{dzyimX} and \emph{tsyijH} (Guangyun, in Baxter's transcription), which glosses the first 'excessive' and the second 'to arrive'.
\end{cousinweb}

\begin{grammarlens}[title={What you've built}]
One more word for giving reasons.
\end{grammarlens}

\subsection*{Your turn}
\begin{itemize}
\raggedright
  \item \emph{Say it:} \emph{shènzhì}
  \item \emph{Say it:} \emph{shènzhì}, once more
  \item \emph{Say it:} \emph{tā hěn máng shènzhì méiyǒu shíjiān chī fàn} — He's so busy he hasn't even got time to eat.
  \item \emph{Recall:} say the Mandarin for the effect something has, then the Mandarin for in short, then say \emph{shènzhì} again
\end{itemize}

\begin{tcolorbox}[breakable,colback=teal!4,colframe=teal!35!black,title={Before you move on}]
What does shènzhì mean? (\textquotedblleft{}Even (going so far as)\textquotedblright{}.) Say it once more.
\end{tcolorbox}
\section[\texorpdfstring{\zh{虽然} (suīrán)}{虽然 (suīrán)}]{\zh{虽然} (suīrán) — although}

\label{lesson:ZH-C371-suiran}



\begin{quote}
Before the new one: say the Mandarin for in short, then the Mandarin for even (going so far as).
\end{quote}

\subsection*{You'll want to know: \zh{虽然}}
\textbf{\zh{虽然}} — \emph{suīrán} — \textquotedblleft{}although\textquotedblright{}.

\begin{quote}
\textbf{\zh{虽然下雨了但是我们还是去了}} \emph{suīrán xiàyǔ le dànshì wǒmen háishi qù le} — Although it rained, we still went.
\end{quote}

The new character, in pieces: \zh{虽} — nine strokes in two pieces: \zh{口} \emph{kǒu} on top, and under it an insect piece, six strokes: a small box with a vertical through it, a stroke rising along the foot, and a dot. The next lesson puts it in your hand.

\begin{cousinweb}
\textbf{\zh{虽然}} (\emph{suīrán}) is built from \emph{suī} and \zh{然}, both inherited from Middle Chinese, \emph{swij} and \emph{nyen} (Guangyun, in Baxter's transcription), which glosses the first as a particle and the second as a particle and 'like this'.
\end{cousinweb}

\begin{grammarlens}[title={What you've built}]
One more word for giving reasons.
\end{grammarlens}

\subsection*{Your turn}
\begin{itemize}
\raggedright
  \item \emph{Say it:} \emph{suīrán}
  \item \emph{Say it:} \emph{suīrán}, once more
  \item \emph{Say it:} \emph{suīrán xiàyǔ le dànshì wǒmen háishi qù le} — Although it rained, we still went.
  \item \emph{Recall:} say the Mandarin for in short, then the Mandarin for even (going so far as), then say \emph{suīrán} again
\end{itemize}

\begin{tcolorbox}[breakable,colback=teal!4,colframe=teal!35!black,title={Before you move on}]
What does \zh{虽然} mean? (\textquotedblleft{}Although\textquotedblright{}.) Say it once more.
\end{tcolorbox}
\section[\texorpdfstring{\zh{虽} (suī)}{虽 (suī)}]{\zh{虽} — a mouth over an insect}

\label{lesson:ZH-W371-sui-though}



\begin{quote}
The word first, then the character inside it.

\begin{itemize}
\raggedright
  \item \emph{Recall:} say \emph{suīrán} — \textbf{R1}, one lesson back, before it decays
\end{itemize}

\begin{quote}
\zh{口} \zh{然}
\end{quote}

\emph{Read it:} each one aloud, with nothing said first — \emph{kǒu}, \emph{rán}
\end{quote}

\subsection*{Script — a mouth over an insect}
\hlblockfigure{\detokenize{figures/ZH-W371-sui-though-filmstrip.pdf}}{How \zh{虽} is written, stroke by stroke}

\begin{quote}
\zh{虽}
\end{quote}

\textbf{Two pieces, top before bottom.} Count the pieces, not the strokes.

\textbf{The top piece, \zh{口}, three strokes}, small: a vertical, a horizontal that turns down, the bottom horizontal.

\textbf{The bottom piece, the insect piece, six strokes}: a vertical, a horizontal that turns down, a horizontal closing the small box, a vertical down through it, a stroke rising to the right along the foot, then a dot on the right.

3 + 6 = nine.

\subsection*{Your turn}
\begin{itemize}
\raggedright
  \item \emph{Write it:} \zh{虽}, all nine, with the model in front of you
  \item \emph{Cover:} write it again from memory, and count the pieces, not the strokes
\end{itemize}

\begin{tcolorbox}[breakable,colback=teal!4,colframe=teal!35!black,title={Before you move on}]
How many pieces in \zh{虽}? (\textbf{Two}.) What is on top? (\textbf{\zh{口}}.) What is the last stroke? (\textbf{A dot on the right}.)

Source: \href{https://raw.githubusercontent.com/chanind/hanzi-writer-data/68d10a4b21150cae5e1ebbd223eed289cf32d90c/data/\%E8\%99\%BD.json}{Hanzi Writer Data \zh{虽}, pinned stroke paths and medians}.
\end{tcolorbox}
\section[jiélùn]{jiélùn — conclusion}

\label{lesson:ZH-C371-jielun}



\begin{quote}
Before the new one: say the Mandarin for even (going so far as), then the Mandarin for although.
\end{quote}

\subsection*{You'll want to know: jiélùn}
\emph{jiélùn} — \textquotedblleft{}conclusion\textquotedblright{}.

\begin{quote}
\emph{wǒ juéde zhège jiélùn bú duì} — I don't think that conclusion is right.
\end{quote}

Its characters come later; for now it is a word for the ear.

\begin{cousinweb}
\emph{jiélùn} is built from \emph{jié} and \emph{lùn}, both inherited from Middle Chinese, \emph{ket} and \emph{lwon} (Guangyun, in Baxter's transcription), which glosses the first 'to tie' and the second 'to discuss'.
\end{cousinweb}

\begin{grammarlens}[title={What you've built}]
One more word for giving reasons.
\end{grammarlens}

\subsection*{Your turn}
\begin{itemize}
\raggedright
  \item \emph{Say it:} \emph{jiélùn}
  \item \emph{Say it:} \emph{jiélùn}, once more
  \item \emph{Say it:} \emph{wǒ juéde zhège jiélùn bú duì} — I don't think that conclusion is right.
  \item \emph{Recall:} say the Mandarin for even (going so far as), then the Mandarin for although, then say \emph{jiélùn} again
\end{itemize}

\begin{tcolorbox}[breakable,colback=teal!4,colframe=teal!35!black,title={Before you move on}]
What does jiélùn mean? (\textquotedblleft{}Conclusion\textquotedblright{}.) Say it once more.
\end{tcolorbox}
\section[zǒngjié]{zǒngjié — to sum up}

\label{lesson:ZH-C371-zongjie}



\begin{quote}
Before the new one: say the Mandarin for although, then the Mandarin for a conclusion.
\end{quote}

\subsection*{You'll want to know: zǒngjié}
\emph{zǒngjié} — \textquotedblleft{}to sum up\textquotedblright{}.

\begin{quote}
\emph{lǎoshī zuìhòu zǒngjié le jīntiān de kè} — At the end the teacher summed up the lesson.
\end{quote}

Its characters come later; for now it is a word for the ear.

\begin{cousinweb}
\emph{zǒngjié} is built from \emph{zǒng} and \emph{jié}, both inherited from Middle Chinese, \emph{tsuwngX} and \emph{ket} (Guangyun, in Baxter's transcription), which glosses the first 'to bind together' and the second 'to tie'.
\end{cousinweb}

\begin{grammarlens}[title={What you've built}]
One more word for giving reasons.
\end{grammarlens}

\subsection*{Your turn}
\begin{itemize}
\raggedright
  \item \emph{Say it:} \emph{zǒngjié}
  \item \emph{Say it:} \emph{zǒngjié}, once more
  \item \emph{Say it:} \emph{lǎoshī zuìhòu zǒngjié le jīntiān de kè} — At the end the teacher summed up the lesson.
  \item \emph{Recall:} say the Mandarin for although, then the Mandarin for a conclusion, then say \emph{zǒngjié} again
\end{itemize}

\begin{tcolorbox}[breakable,colback=teal!4,colframe=teal!35!black,title={Before you move on}]
What does zǒngjié mean? (\textquotedblleft{}To sum up\textquotedblright{}.) Say it once more.
\end{tcolorbox}
\section[yuánzé]{yuánzé — principle}

\label{lesson:ZH-C371-yuanze}



\begin{quote}
Before the new one: say the Mandarin for a conclusion, then the Mandarin for to sum up.
\end{quote}

\subsection*{You'll want to know: yuánzé}
\emph{yuánzé} — \textquotedblleft{}principle\textquotedblright{}.

\begin{quote}
\emph{wǒ juéde měi ge rén dōu yīnggāi yǒu yuánzé} — I think everyone should have principles.
\end{quote}

Its characters come later; for now it is a word for the ear.

\begin{cousinweb}
\emph{yuánzé} is built from \emph{yuán} and \emph{zé}, both inherited from Middle Chinese, \emph{ngjwon} and \emph{tsok} (Guangyun, in Baxter's transcription), which glosses the first 'a broad plain' and the second 'a rule'. How the first came to mean a source is not settled.
\end{cousinweb}

\begin{grammarlens}[title={What you've built}]
One more word for giving reasons.
\end{grammarlens}

\subsection*{Your turn}
\begin{itemize}
\raggedright
  \item \emph{Say it:} \emph{yuánzé}
  \item \emph{Say it:} \emph{yuánzé}, once more
  \item \emph{Say it:} \emph{wǒ juéde měi ge rén dōu yīnggāi yǒu yuánzé} — I think everyone should have principles.
  \item \emph{Recall:} say the Mandarin for a conclusion, then the Mandarin for to sum up, then say \emph{yuánzé} again
\end{itemize}

\begin{tcolorbox}[breakable,colback=teal!4,colframe=teal!35!black,title={Before you move on}]
What does yuánzé mean? (\textquotedblleft{}Principle\textquotedblright{}.) Say it once more.
\end{tcolorbox}
